\chapter{Final Results, Recovery Provenance, and Statistical Summary}

\section{Canonical Evidence Base}

The final empirical evidence base is the frozen \texttt{final\_113} artifact set and its canonical reporting outputs. That evidence base contains 113 completed rows, 16 recovered replacements, 12 corrected guarded hierarchical reruns from 14 March 2026, and 4 recovered DQN reruns retained as descriptive ablations. This chapter therefore reports completed results rather than a provisional campaign status.

\input{tables/generated/current_results_status.tex}

Two provenance rules are essential for interpreting the chapter correctly. First, recovered replacement rows are treated as part of the canonical evidence set only because they are explicitly tracked and rebuilt into the frozen reporting layer. Second, the hierarchical reruns are reported as a corrected guarded branch rather than as if they were identical to the earlier raw hierarchical runs. That distinction matters because the guarded reruns enforce calendar windows, seasonal fertilizer gating, and annual N/P/K budgets.

\begin{figure}[H]
    \centering
    \includegraphics[width=0.78\textwidth]{final_artifact_audit.pdf}
    \caption{Audit counts that remain relevant after canonicalization. The main residual provenance issue is that recovered guarded hierarchical reruns do not include the original thesis-report directories.}
    \label{fig:final-artifact-audit}
\end{figure}

\section{Fertilization Results}

The fertilization branch is the strongest and most statistically informative part of the thesis. Across repeated three-seed groups, the best mean deterministic return came from \emph{A2C | nonadaptive | fixed weather | years=5000}, with a mean deterministic return of 779,267.4 and a 95\% confidence interval of [727,686.5, 830,848.2]. The best single fertilization run was \emph{PPO | adaptive | fixed weather | years=3000 | seed=1}, which reached a deterministic return of 791,255.5.

\begin{figure}[H]
    \centering
    \includegraphics[width=0.96\textwidth]{final_fertilization_scores.pdf}
    \caption{Repeated fertilization groups ranked by deterministic return. Red markers show the mean Pakistan holdout return used as the main robustness metric for this branch.}
    \label{fig:final-fertilization}
\end{figure}

\input{tables/generated/final_fertilization_table.tex}

The inferential results for fertilization are meaningful rather than cosmetic. Type II ANOVA indicates strong method and budget effects, smaller but still significant adaptivity and weather effects, and a method-by-budget interaction. At the same time, the top targeted pairwise comparisons do not all survive Holm correction, so the interpretation should remain at the level of broad factor effects rather than exaggerated winner-take-all claims.

\section{Non-Hierarchical Crop-Planning Results}

For the non-hierarchical crop-planning branch, the headline metric is \texttt{eval\_det/mean\_reward}. On that metric the best repeated group was \emph{PPO | nonadaptive | fixed weather}, with a mean of 21,230.9 and a 95\% confidence interval of [15,225.6, 27,236.2]. The best single non-hierarchical crop-planning run was \emph{A2C | adaptive | fixed weather | seed=2}, with \texttt{eval\_det/mean\_reward} equal to 23,601.6.

\begin{figure}[H]
    \centering
    \includegraphics[width=0.92\textwidth]{final_crop_scores.pdf}
    \caption{Repeated non-hierarchical crop-planning groups ranked by the headline metric \texttt{eval\_det/mean\_reward}. Deterministic return is retained only as a supporting metric for this branch.}
    \label{fig:final-crop}
\end{figure}

\input{tables/generated/final_crop_table.tex}

The crop-planning branch is more tightly clustered than fertilization. The repeated-group ranking is clear descriptively, but the ANOVA and targeted pairwise tests show substantial overlap across the main factors at this sample size. Accordingly, the thesis can claim that the branch was executed and compared consistently, while keeping the main emphasis on the completed comparison framework rather than on small numerical gaps among the top repeated groups.

\section{Corrected Guarded Hierarchical Reruns}

The hierarchical branch is reported separately because it runs under a corrected reward and constraint regime and therefore is not ranked head-to-head inside the non-hierarchical crop-planning leaderboard. The guarded reruns comprise 12 rows across four repeated groups. The best repeated guarded group was \emph{PPO | fixed weather | guarded rerun}, with a mean deterministic return of 1,861,223.1 and a 95\% confidence interval of [1,512,430.8, 2,210,015.4]. The best single guarded rerun reached 1,971,988.3 deterministic return.

\begin{figure}[H]
    \centering
    \includegraphics[width=0.90\textwidth]{final_hierarchical_scores.pdf}
    \caption{Corrected guarded hierarchical reruns ranked by deterministic return. The baseline comparator is shown separately because this branch is reported as its own corrected regime.}
    \label{fig:final-hierarchical}
\end{figure}

\input{tables/generated/final_hierarchical_table.tex}

These reruns are authoritative for the current repo and should be treated as a dedicated guarded branch. Two provenance notes remain important: the recovered bundles do not contain the original thesis-report directories, and the guarded branch remains a corrected rerun regime rather than a direct continuation of the earlier raw hierarchical narrative. The committee can therefore be told that hierarchical work was completed, corrected, frozen, and kept fully auditable through explicit provenance.

\section{Statistical Summary}

The full inferential layer is restricted to repeated three-seed groups. This includes 24 fertilization groups, 8 non-hierarchical crop-planning groups, and 4 hierarchical guarded-rerun groups. The single baseline row and all DQN reruns remain descriptive only.

\input{tables/generated/statistical_summary.tex}

The main statistical message is simple. Fertilization shows defensible factor-level effects, especially for algorithm family and budget. Non-hierarchical crop planning and the corrected guarded hierarchical reruns show closer overlap at the current sample size. That branch-to-branch difference is acceptable because it is reported explicitly rather than hidden.

\section{Artifact and Evidence Notes}

Canonicalization removes document inconsistency, but it does not erase the remaining evidence limits. The thesis still depends on simulation rather than field validation, the crop-planning scope remains maize-soy, irrigation is not a learned action, and the guarded hierarchical reruns still have incomplete thesis-report directories. These caveats are surfaced directly rather than being treated as afterthoughts.

\input{tables/generated/artifact_validity_caveats.tex}

\section{Result-Level Conclusion}

The completed campaign supports three defensible claims. First, the fertilization branch is the strongest empirical contribution and carries the clearest quantitative signal. Second, the non-hierarchical crop-planning branch is completed and comparable, with a clear headline ranking on the chosen metric. Third, the hierarchical branch is now reported as a corrected guarded rerun branch with explicit provenance and dedicated interpretation. That combination of completion, provenance control, and bounded interpretation is the main strength of the final results chapter.
